\paragraph{\trSuspension。}\label{sub:shift}
\trTensorProduct{}最常用的例子，是取 $Z$ 为一维\trVectorSpace{} $\KK[\sigma]$，
并置于某个\trDegree{} $-\sigma$ 上。其结果记作
$\KK[\sigma]\otimes\scr{M}=SS^{\sigma}\scr{M}$，相应地
$\KK[\sigma]\otimes Y=S^{\sigma}Y$，并称为 $\scr{M}$ 或 $Y$ 的%
\textbf{$\sigma$ 重\trSuspension}（$\sigma$-fold shift）。当 $\sigma=1$ 时我们
省去\trIndex，径直写作 $SS\scr{M}$、$SY$。这些\trSuspension（如果存在的话）带有%
\trCanonical{}的\trIsomorphism：
\begin{equation}\label{eq:shift-hom-isomorphism}
\Hom_{H(\scr{A})}(Y_0,S^{\sigma}Y_1)\cong\Hom_{H(\scr{A})}(Y_0,Y_1)[\sigma],
\end{equation}
\begin{equation}\label{eq:shift-functorial-isomorphism}
\Hom_{H(\scr{A})}(Y_0,Y_1)\cong\Hom_{H(\scr{A})}(S^{\sigma}Y_0,S^{\sigma}Y_1),
\end{equation}
\begin{equation}\label{eq:shift-inverse-isomorphism}
\Hom_{H(\scr{A})}(S^{\sigma}Y_0,Y_1)[\sigma]\cong\Hom_{H(\scr{A})}(Y_0,Y_1).
\end{equation}
第一个按定义即得；第二个来自 $SS^{\sigma}=\KK[\sigma]\otimes-$ 的\trFunctor{}本性。
第三个的得到方式是：先把\eqref{eq:shift-hom-isomorphism} 用于对
$(S^{\sigma}Y_0,Y_1)$，再把\eqref{eq:shift-functorial-isomorphism} 反向地
（从右到左）用于 $(Y_0,Y_1)$。

\begin{remark}\label{thm:shift-sign-caution}
看待\eqref{eq:shift-inverse-isomorphism} 的另一种方式是：从
$(S^{\sigma}Y_0,Y_1)$ 出发，用\eqref{eq:shift-functorial-isomorphism} 过渡到
$(Y_0,S^{-\sigma}Y_1)$，再反向使用\eqref{eq:shift-hom-isomorphism} 到达
$(Y_0,Y_1)$。可惜这与前面的定义相差一个正负号 $(-1)^{\sigma}$，因此在使用这些%
\trIsomorphism{}时需要多加小心。
\end{remark}

假设 \trACategory{} $\scr{A}$ 具有如下性质：对一切 $Y$，$SY$ 都存在。此时便可
构造一个\trSuspension{} \trAFunctor{} $S\colon\scr{A}\to\scr{A}$，它在
$H^0(\mathit{fun}(\scr{A},\scr{A}))$ 中至多相差一个\trIsomorphism{}的意义下
唯一，且使得图表
\begin{tikzcd}
\scr{A}\ar[r,"\scr{J}"]\ar[d,"S"'] & \scr{Q}\ar[d,"SS"]\\
\scr{A}\ar[r,"\scr{J}"'] & \scr{Q}
\end{tikzcd}
在 $H^0(\mathit{fun}(\scr{A},\scr{Q}))$ 中交换到相差一个\trIsomorphism。具体
说来，考虑由全部形如 $SY$ 的\trObject{}组成的\trFull{} \trAinfty-子\trCategory{}%
$\widetilde{\scr{A}}\subset\scr{A}$，并相应地取子\trCategory{}%
$\widetilde{\scr{Q}}\subset\scr{Q}$，它由那些在 $H^0(\scr{Q})$ 中\trIsomorphic{}于
某个 $SS\scr{Y}=SS\scr{J}(Y)$ 的 \trAinfty-模组成。由构造，\trYonedaEmbedding{}%
的限制是一个\trQuasiequivalence{} $\widetilde{\scr{A}}\to\widetilde{\scr{Q}}$，
因而可按\cref{thm:quasi-equivalence-homotopy-inverse}求逆。取定这样一个逆 $\scr{K}$ 之后，令
$S=\scr{K}\circ SS\circ\scr{J}$。以上就\trSuspension{}的情形作了解释，但同样的
论证也适用于 \trAinfty-模上其他\trFunctor{}性的构造。

\begin{remark}\label{thm:tensor-product-existence}
任意\trSuspension{} $S^{\sigma}Y$（$\sigma\in\ZZ$）的存在性，连同\trDirectSum{}%
$Y_0\oplus Y_1$ 的存在性，蕴含：对具有有限维\trCohomology{}的\trComplex{}%
$(Z,\partial_Z)$，\trTensorProduct{} $Z\otimes Y$ 也存在。具体说来，若取
$\widetilde{Z}=H(Z)$ 的一组\trBasis{} $\{[z^i]\}$，其中 $\sigma^i=-|z^i|$，则
\begin{equation*}
\bigoplus_{i\in I}S^{\sigma^i}Y
\end{equation*}
在\trYonedaEmbedding{}下被映为
$\bigoplus_{i\in I}SS^{\sigma^i}\scr{Y}\cong\widetilde{Z}\otimes\scr{Y}$。由%
\cref{thm:tensor-product-homology}，它也是 $Z\otimes\scr{Y}$ 的一个%
\trQuasiRepresent\trObject。
\end{remark}
